\subsection{Free observation spaces}
\label{obs:free-spaces}

The construction in this subsection is purely set-theoretic.  The typing
conditions on frames will be imposed by evaluation; they do not restrict
the tables constructed here.  Fix the set $\mathcal Q$ of normal syntax
over $\Omega$, with bound names identified up to $\alpha$-conversion.
Let $\mathcal T_{\mathrm{pl}}$ be the set of actual plain active types,
and suppose that the nonempty spaces $\mathcal K_D$ have already been
constructed for $D\in\mathcal T_{\mathrm{pl}}$.  The fixed set of tagged
plain values is
\begin{equation}
  \mathcal V_{\mathrm{pl}}
   =\coprod_{D\in\mathcal T_{\mathrm{pl}}}
       \bigl(\{D\}\times\mathcal K_D\bigr).
  \label{obs:free-plain-values}
\end{equation}
Thus a member of $\mathcal V_{\mathrm{pl}}$ remembers both its type $D$ and its value
$v\in\mathcal K_D$.  Let $\mathcal B$ be the set of base observations:
a member chooses a candidate at every actual type.  Write
$\mathbf t\in\mathcal B$ for the choice of the top candidate everywhere.
For the present construction only the set $\mathcal B$ and its
distinguished member $\mathbf t$ are needed.  Put
\[
  \mathcal O=\{\star\}\sqcup\mathcal V_{\mathrm{pl}},
\]
where the union is disjoint.  In particular, both output sets
$\mathcal B$ and $\mathcal O$ have specified default elements, even if
there are no plain active types.

\paragraph{Frames, stacks, and tables.}
For any set $Y$ define the payload set
\[
  V(Y)
   =\{\mathrm{none}\}
      \sqcup\bigl(\{\mathrm{opp}\}\times Y\bigr)
      \sqcup\bigl(\{\mathrm{plain}\}\times\mathcal V_{\mathrm{pl}}\bigr).
\]
The three tags make these alternatives disjoint.  Define continue and
exit frames by
\[
  F_c(Y)=\mathcal Q\times V(Y)\times\{c\},
  \qquad
  F_e(Y)=\mathcal Q\times V(Y)\times\{e\}.
\]
Here $Y$ is the space of opposite-sign payloads.  The set of non-exiting
stacks and the set of exiting stacks are, respectively,
\begin{equation}
  S_c(Y)=\coprod_{n\geq 0} F_c(Y)^n,
  \qquad
  S_e(Y)=\coprod_{n\geq 0}
                 \bigl(F_c(Y)^n\times F_e(Y)\bigr).
  \label{obs:free-stacks}
\end{equation}
Thus $S_c(Y)$ contains the empty stack $\epsilon$, and an exit frame
occurs exactly once, at the end of a stack in $S_e(Y)$.  Set
\begin{equation}
  \operatorname{Tab}(Y)
      =\mathcal B^{S_c(Y)}\times\mathcal O^{S_e(Y)}.
  \label{obs:free-tab}
\end{equation}
A table $t=(t_c,t_e)$ gives a base observation on each non-exiting
stack and an exit observation on each exiting stack.  Its default
$d_Y$ is specified by
\[
  (d_Y)_c(\sigma)=\mathbf t,
  \qquad
  (d_Y)_e(\tau)=\star.
\]
Consequently $\operatorname{Tab}(Y)$ is nonempty for every $Y$,
including $Y=\emptyset$.

We seek two spaces whose tables take values of the opposite sign as
payloads.  Their construction must solve two different problems.  A plus
value should accept every finite stack of minus values.  A minus value
must be representable at a finite stage, even though it is tested on
arbitrarily many plus values and arbitrarily long stacks.  Choosing a
stage separately for each stack would not give that representation.

The construction below supplies a space $C_*$ of minus values and the
full space $P_*=\operatorname{Tab}(C_*)$ of plus tables.  Increasing
stages $C_r$ give restrictions
$\rho_r:P_*\to P_r=\operatorname{Tab}(C_r)$ of plus tables.
A table on plus payloads represents a minus value precisely when
one index $r$ works for every observation: replacing each plus payload
$p$ by $p'$ with $\rho_r(p)=\rho_r(p')$ leaves the table output
unchanged.  Finally, fixing any continue prefix must again give a value
of the same sign; on the minus side it must retain the same representing
stage.  These are the representation and slicing properties that typed
evaluation will use.  We first construct the stages and then prove this
characterization.

\paragraph{Relabeling and contravariance.}
A map $f:Y\to Z$ induces stack maps
$S_a(f):S_a(Y)\to S_a(Z)$ for $a\in\{c,e\}$.
These replace $(\mathrm{opp},y)$ by $(\mathrm{opp},f(y))$ in each
payload position.  Keys, flags, and the other payload alternatives
are unchanged.  Define
\begin{equation}
  \operatorname{Tab}(f):\operatorname{Tab}(Z)
                  \longrightarrow\operatorname{Tab}(Y),
  \qquad
  \bigl(\operatorname{Tab}(f)(t)\bigr)_a
                   =t_a\circ S_a(f).
  \label{obs:free-contravariance}
\end{equation}
Relabeling preserves identities and composition.  Hence
\[
  \operatorname{Tab}(\operatorname{id}_Y)
     =\operatorname{id}_{\operatorname{Tab}(Y)},
  \qquad
  \operatorname{Tab}(g\circ f)
     =\operatorname{Tab}(f)\circ\operatorname{Tab}(g).
\]
It also preserves the default tables:
$\operatorname{Tab}(f)(d_Z)=d_Y$.

We will use the following two elementary facts.
\begin{enumerate}
\item If $f$ is injective, then $S_c(f)$ and $S_e(f)$ are injective,
and $\operatorname{Tab}(f)$ is surjective.  Indeed, given a table on
$Y$, define a table on $Z$ by transporting its values to the images
of $S_c(f)$ and $S_e(f)$, and assigning $\mathbf t$ or $\star$ on the
respective complements.  Injectivity makes the transported values
unambiguous, and restriction recovers the given table.
\item If $f$ is surjective, then $S_c(f)$ and $S_e(f)$ are surjective,
and $\operatorname{Tab}(f)$ is injective.  To lift a stack, lift its
finitely many opposite payloads through $f$.  If two tables become
equal after precomposition with these surjective stack maps, they
agree on every stack and therefore are equal.
\end{enumerate}
In particular these statements apply when an opposite payload set is
empty; the stacks without opposite payloads remain available.

\paragraph{The alternating chain.}
Define, for $r=0,1,2,\ldots$,
\begin{equation}
  C_0=\emptyset,
  \qquad P_r=\operatorname{Tab}(C_r),
  \qquad C_{r+1}=\operatorname{Tab}(P_r).
  \label{obs:free-stages}
\end{equation}
Let $i_0:C_0\to C_1$ be the empty map.  Inductively, once
$i_r:C_r\to C_{r+1}$ is defined, put
\begin{equation}
  j_r=\operatorname{Tab}(i_r):P_{r+1}\longrightarrow P_r,
  \qquad
  i_{r+1}=\operatorname{Tab}(j_r):C_{r+1}\longrightarrow C_{r+2}.
  \label{obs:free-bonding}
\end{equation}
The preceding facts prove by induction that every $i_r$ is injective
and every $j_r$ is surjective.  Each $j_r$ is restriction to stacks
whose opposite payloads lie in the embedded copy of $C_r$; its
surjectivity has the explicit default-extension proof above.

For $r\leq s$, write $i_{r,s}:C_r\to C_s$ for the composite of the
injections, with $i_{r,r}$ the identity.  Let $C_*$ be the set colimit
of this chain.  Explicitly, take the disjoint union of the $C_r$ and
identify $(r,x)$ with $(s,y)$ precisely when, at some common stage
$m\geq r,s$,
\[
  i_{r,m}(x)=i_{s,m}(y).
\]
The canonical maps $e_r:C_r\to C_*$ are injective and satisfy
$e_{r+1}\circ i_r=e_r$.  Their images form an increasing union
equal to $C_*$.  We may therefore regard the stages as nested
subsets after these identifications.  Every finite family of elements
of $C_*$ lies in a common stage: choose a representing stage for
each element and then take their maximum.  Every element in fact
has a representative in some $C_{r+1}$, since $C_0$ is empty.

\paragraph{Final plus tables and finite-stage agreement.}
Put
\begin{equation}
  P_*=\operatorname{Tab}(C_*),
  \qquad
  \rho_r=\operatorname{Tab}(e_r):P_*\longrightarrow P_r.
  \label{obs:free-final-plus}
\end{equation}
Each $\rho_r$ is surjective, by default extension along the injection
$e_r$.  Contravariance and the identities for $e_r$ give
\begin{equation}
  \rho_r=j_r\circ\rho_{r+1}.
  \label{obs:free-restriction-compatibility}
\end{equation}
For $p,p'\in P_*$ write
\begin{equation}
  p\equiv_r p' \quad\Longleftrightarrow\quad
           \rho_r(p)=\rho_r(p').
  \label{obs:free-equivalence}
\end{equation}
Agreement at a larger stage implies agreement at every smaller
stage, by repeated use of
\eqref{obs:free-restriction-compatibility}.

Every stack in $S_a(C_*)$, for $a\in\{c,e\}$, belongs to the image
of $S_a(e_r)$ for some $r$: it has only finitely many opposite
payloads, which lie in a common stage.  The same holds for any
finite family of stacks, with one common $r$.  It follows that
\begin{equation}
  \bigl(p\equiv_r p'\text{ for every }r\bigr)
       \quad\Longleftrightarrow\quad p=p'.
  \label{obs:free-separation}
\end{equation}
For the forward implication, lift any specified final stack from a
common stage and evaluate the equality $\rho_r(p)=\rho_r(p')$ on
that lift.  The reverse implication is immediate.  No assertion
that a plus table itself belongs to a finite stage is needed.

\paragraph{Final minus behaviors.}
For each $r\geq0$, define
\begin{equation}
  \beta_r=\operatorname{Tab}(\rho_r):
        C_{r+1}=\operatorname{Tab}(P_r)
            \longrightarrow\operatorname{Tab}(P_*).
  \label{obs:free-behavior-stage}
\end{equation}
In coordinates, for $u\in C_{r+1}$ and $\sigma\in S_a(P_*)$,
\[
  \bigl(\beta_r(u)\bigr)_a(\sigma)
       =u_a\bigl(S_a(\rho_r)(\sigma)\bigr).
\]
Thus each plus payload of a final stack is first restricted to
$P_r$.  Since $\rho_r$ is surjective, $\beta_r$ is injective.
Moreover,
\begin{align*}
  \beta_{r+1}\circ i_{r+1}
    &=\operatorname{Tab}(\rho_{r+1})
                          \circ\operatorname{Tab}(j_r)\\
    &=\operatorname{Tab}(j_r\circ\rho_{r+1})
      =\operatorname{Tab}(\rho_r)=\beta_r.
\end{align*}
These maps consequently induce a well-defined map
\begin{equation}
  \beta:C_*\longrightarrow\operatorname{Tab}(P_*),
  \qquad
  \beta(e_{r+1}(u))=\beta_r(u).
  \label{obs:free-behavior}
\end{equation}
It is injective: represent any two proposed equal-behavior elements
at a common stage $C_{r+1}$ and use the injectivity of $\beta_r$.
We henceforth evaluate a minus value through this behavior map.

The defaults are compatible with all bonding maps.  In particular,
the default elements $d_{P_r}\in C_{r+1}$ represent one element of
$C_*$, whose behavior is $d_{P_*}$.  The plus default is
$d_{C_*}\in P_*$.  Both free spaces are therefore nonempty.

\paragraph{Exact characterization by a uniform stage bound.}
Fix $r$.  For $a\in\{c,e\}$ and $\sigma,\tau\in S_a(P_*)$ write
$\sigma\sim_r\tau$ when they have identical lengths, keys, flags,
and payload alternatives, their plain payloads are identical, and
each corresponding pair of opposite payloads $p,p'$ satisfies
$p\equiv_r p'$.  Equivalently,
\begin{equation}
  \sigma\sim_r\tau
    \quad\Longleftrightarrow\quad
  S_a(\rho_r)(\sigma)=S_a(\rho_r)(\tau).
  \label{obs:free-stack-equivalence}
\end{equation}
This equivalence follows directly from the disjoint payload tags
and the definition of stack relabeling.

For a final table $t\in\operatorname{Tab}(P_*)$, the following
conditions are equivalent:
\begin{enumerate}
\item $t$ belongs to the image of $\beta_r$.
\item For both $a=c$ and $a=e$, the equality
      $\sigma\sim_r\tau$ implies $t_a(\sigma)=t_a(\tau)$.
\end{enumerate}
The first condition implies the second by
\eqref{obs:free-behavior-stage}.  Conversely, each
$S_a(\rho_r)$ is surjective, since a finite stack of elements of
$P_r$ can be lifted through the surjection $\rho_r$.  For
$\eta\in S_a(P_r)$ define $u_a(\eta)$ to be $t_a(\sigma)$ for
any lift $\sigma$ of $\eta$.  Such a lift exists, and any two lifts
are $\sim_r$-equivalent by
\eqref{obs:free-stack-equivalence}.  The second condition therefore
makes the definition unambiguous.  It defines a unique
$u\in\operatorname{Tab}(P_r)=C_{r+1}$ with $\beta_r(u)=t$.
This proves the equivalence and uniqueness of the factor.

Consequently
\begin{equation}
  \beta(C_*)
    =\bigcup_{r\geq0}\beta_r(C_{r+1})
    =\left\{
       t\in\operatorname{Tab}(P_*):
       \begin{array}{l}
       \text{for some }r,\text{ each }t_a\text{ is constant}\\
       \text{on every }\sim_r\text{-class},\quad a\in\{c,e\}
       \end{array}
       \right\}.
  \label{obs:free-uniform-factorization}
\end{equation}
The one stage $r$ here is uniform over all stack lengths, keys,
payloads, and both kinds of ending.  The preceding argument proves
factorization for precisely this uniform condition.

\paragraph{Continue-prefix slicing.}
If $\sigma\in S_c(Y)$ and $t\in\operatorname{Tab}(Y)$, define
the slice $t|\sigma\in\operatorname{Tab}(Y)$ by
\begin{equation}
  (t|\sigma)_a(\tau)=t_a(\sigma\tau),
       \qquad \tau\in S_a(Y),\quad a\in\{c,e\},
  \label{obs:free-slice}
\end{equation}
where juxtaposition denotes concatenation.  The prefix has only
continue frames, so concatenation preserves the kind of ending.
This immediately gives slices of plus tables, by taking $Y=C_*$.

For minus tables, slicing preserves the representing stage.
For $u\in C_{r+1}$ and $\sigma\in S_c(P_*)$,
stack relabeling commutes with concatenation and gives
\begin{equation}
  \beta_r(u)|\sigma
     =\beta_r\bigl(u|S_c(\rho_r)(\sigma)\bigr).
  \label{obs:free-slice-bound}
\end{equation}
To check this equality, evaluate either side on
$\tau\in S_a(P_*)$; both values are
\[
  u_a\bigl(S_c(\rho_r)(\sigma)\,
                 S_a(\rho_r)(\tau)\bigr).
\]
The sliced table inside $\beta_r$ belongs to
$\operatorname{Tab}(P_r)=C_{r+1}$.  Thus a minus value represented
at $C_{r+1}$ has every continue-prefix slice represented at that
same stage.  Because $\beta$ is injective, this defines an intrinsic
slice on $C_*$, independent of the chosen representing stage.
Equivalently, fixing a prefix preserves invariance under
$\sim_r$ on all remaining payloads.  Default tables slice to
default tables.  Define
\[
 \mathcal K_T=\begin{cases}P_*,&T\text{ is free plus},\\
 C_*,&T\text{ is free minus}.\end{cases}
\]
Evaluate a minus value through the injective behavior map $\beta$.
Thus a plus value is an arbitrary table, whereas a proposed minus
behavior defines a value only after one uniform stage bound has been
proved.  The evaluation of a lambda at a free minus type will be the
place where this distinction matters.
